\subsection{The Strategy 2 transfer rule}

\begin{definition}[Raw and center-free transfers]
\label{def:body-transfer-alphabet}
For each $c\in[0,1]$, define $g_c:[0,1]\to[0,1]$ by
\[
 g_c(x)=
 \max\left\{y\in[0,1]:(1-x,y,c)\in\mathcal A\right\}.
\]
For incoming reach $a$, define the center-free lower transfer directly by
\[
 g_c^\vee(a)=1-g_c(1-a).
\]
Thus $x=1-a$ is an incoming boundary-defect coordinate and $g_c(x)$ is the
extremal admissible outgoing demand.
\end{definition}

The midpoint $\frac12(u+v)$ gives the exact threshold
\begin{equation}
 c\ge\frac12\quad\Longrightarrow\quad g_c(x)\le x,
 \qquad
 g_{1/2}(x)=x.
 \label{eq:body-midpoint-threshold}
\end{equation}
Accordingly, a known nonsupercritical role uses the direct bound $B\le1-A$
when its selected radial demand is at most $1/2$, and the raw map $g_c$ when
that demand is larger than $1/2$.  At zero radial demand,
$g_0(x)>x$ for $0<x<1$, so the raw map cannot supply the low-radial identity
handoff.

\begin{definition}[Residual and center-assisted transfers]
\label{def:body-center-assisted-transfer}
Let $J$ be either empty or a closed interval $[L,U]$ with
$0\le L\le U\le1$.  For $p\in[0,1]$, define
\[
 e_J(p)=
 \sup\left\{t\in[0,1]:[0,t]\subseteq[0,p]\cup J\right\},
 \qquad
 \mathcal R_J(p)=1-e_J(p).
\]
Equivalently,
\[
 \mathcal R_\varnothing(p)=1-p,
\qquad
 \mathcal R_{[L,U]}(p)=
 \begin{cases}
 1-p,&p<L,\\
 1-\max\left\{p,U\right\},&p\ge L.
 \end{cases}
\]
For $a,c\in[0,1]$, define the raw center-assisted transfer
\[
 g_{c,J}^\vee(a)=\mathcal R_J(g_c(1-a)).
\]
For a nonsupercritical role with $c\le1/2$, the corresponding handoff is
written directly as $\mathcal R_J(1-a)$.
\end{definition}

\begin{figure}[H]
\centering
\resizebox{.92\linewidth}{!}{\begin{tikzpicture}[x=1cm,y=1cm,>=Latex]
  \node[panellabel,anchor=west] at (.03,3.05)
    {(a) raw transfer and complement dual on the skeleton};

  \begin{scope}[shift={(2.05,.42)},scale=1.65]
    \HexCoords
    \DrawHexagon
    \DrawRays

    % The paper fixes counterclockwise indexing.  At V_0 the incoming edge is
    % e_{5,0}, the outgoing edge is e_{0,1}, and the own radial arm is r_0.
    \coordinate (Amark) at ($(V0)!.38!(V5)$);
    \coordinate (Ymark) at ($(V0)!.622!(V1)$);
    \coordinate (Cmark) at ($(V0)!.44!(O)$);

    \draw[demand] (V0)--(Amark);
    \draw[dimline]
      ($(Amark)+(.075,-.045)$)--($(V5)+(.075,-.045)$)
      node[midway,below right,figlabel,text=black!60] {$x=1-a$};
    \node[figlabel,below right,text=DemandRed]
      at ($(V0)!.28!(V5)$) {$a$};

    \draw[draw=GapPurple,line width=1.35pt] (V0)--(Ymark);
    \node[figlabel,above right,text=GapPurple] at ($(Ymark)+(0,.08)$)
      {$y=g_c(x)$};

    \draw[radialdemand] (V1)--(Ymark);
    \node[figlabel,above left,align=right,text=RadialTeal]
      at ($(V1)+(-.03,.17)$) {$1-y=g_c^\vee(a)$};

    \draw[radialdemand] (V0)--(Cmark);
    \node[figlabel,above,text=RadialTeal]
      at ($(V0)!.50!(O)$) {$c$};

    \fill (V0) circle (.9pt);
    \node[figlabel,right] at (V0) {$V_0$};
    \node[figlabel,above] at (V1) {$V_1$};
    \node[figlabel,below] at (V5) {$V_5=V_{-1}$};
  \end{scope}

  \node[figlabel,atlasbox,align=left,anchor=west] at (5.85,1.30)
    {\textcolor{DemandRed}{red}: incoming reach\\
     \textcolor{GapPurple}{purple}: raw extremal output, not an actual trace\\
     \textcolor{RadialTeal}{teal}: radial parameter or complement dual\\
     \textcolor{black!60}{gray}: complementary defect\\
     \(e_{0,1}\) is center-free for the displayed complement dual\\
     representative positions; displayed identities are exact};
\end{tikzpicture}
}
\caption{Incoming reach and defect, radial parameter, raw extremal output,
and the center-free lower transfer.}
\label{fig:transfer-vtriangle-skeleton}
\end{figure}

\begin{figure}[H]
\centering
\resizebox{.99\linewidth}{!}{\begin{tikzpicture}[x=1cm,y=1cm,>=Latex]
  \node[panellabel,anchor=west] at (.02,2.30)
    {(c) the center-assisted residual on a skeleton edge};

  \foreach \sx/\pt/\et/\case/\formula in {
    1.85/.25/.25/{p<L}/{1-p},
    6.00/.58/.72/{L\le p<U}/{1-U},
    10.15/.86/.86/{p\ge U}/{1-p}
  }{
    \begin{scope}[shift={(\sx,.65)},scale=1.08]
      \HexCoords
      \DrawHexagon
      \DrawRays
      \coordinate (Pp) at ($(V0)!\pt!(V1)$);
      \coordinate (PL) at ($(V0)!.42!(V1)$);
      \coordinate (PU) at ($(V0)!.72!(V1)$);
      \coordinate (Pe) at ($(V0)!\et!(V1)$);

      \draw[black!28,line width=2.8pt] (V0)--(Pe);
      \draw[demand] (V0)--(Pp);
      \draw[centertrace] (PL)--(PU);
      \draw[radialdemand]
        ($(V1)+(.05,.03)$)--($(Pe)+(.05,.03)$);
      \fill[DemandRed] (Pp) circle (.75pt);
      \fill[black!65] (Pe) circle (.7pt);

      \node[figlabel,above right] at ($(V0)!.52!(Pp)$) {$p$};
      \node[figlabel,above] at ($(PL)!.50!(PU)$) {$J=[L,U]$};
      \node[figlabel,anchor=east,align=right] at ($(V1)+(-.06,.25)$)
        {$A_{\rm next}\ge\mathcal R_J(p)$};
      \node[figlabel,below] at (0,-1.55)
        {$\case:\quad\mathcal R_J(p)=\formula$};
    \end{scope}
  }

  \node[figlabel,align=center] at (6.00,-1.78)
    {Gray is the component of \([0,p]\cup J\) attached to \(V_0\);
     the teal far-side reach is forced from \(V_1\).\\
     For \(p=g_c(1-a)\) this gives \(g_{c,J}^{\vee}(a)\);
     for \(p=\widehat g_c(1-a)\) it gives
     \(\widehat g_{c,J}^{\vee}(a)\).};
\end{tikzpicture}
}
\caption{The three geometric cases in the residual operator
$\mathcal R_{[L,U]}(p)$, drawn on a handoff edge.}
\label{fig:transfer-residual-skeleton}
\end{figure}

\begin{definition}[Free strict-supercritical envelope]
\label{def:body-supercritical-envelope}
For $0\le c<1/2$, define
\[
 g_c^{\rm sc}
 =\sup_{\substack{x\in[0,1]\\g_c(x)>x}}g_c(x)
 =\frac{c+\sqrt{c^2-8c+4}}2.
\]
\end{definition}

Figures~\ref{fig:transfer-vtriangle-skeleton}
and~\ref{fig:transfer-residual-skeleton} illustrate the transfer and residual
coordinates.  The variable $x$ is complementary reach data, not necessarily
a geometrically uncovered segment, and $g_c(x)$ is an extremal admissible
output rather than the actual reach $B(T)$.

\begin{proposition}[Strategy 2 transfer bounds]
\label{prop:body-transfer-interface}
Fix $a,c\in[0,1]$, and let $J$ be either empty or a closed interval in
$[0,1]$.  Let $T$ be an actual V triangle satisfying $A(T)\ge a$ and
$C(T)\ge c$, and let $T_{\mathrm{next}}$ be the next incident actual V
triangle.  When a center trace assists the handoff, $J$ is its scalar
parameter interval on that edge.
\begin{enumerate}
\item Always,
\[
 B(T)\le g_c(1-a).
\]
If the outgoing edge is center-free, then
\[
 A(T_{\mathrm{next}})\ge g_c^\vee(a).
\]
\item If $T$ is nonsupercritical, then
\[
 B(T)\le
 \begin{cases}
 1-a,&0\le c\le1/2,\\
 g_c(1-a),&1/2<c\le1.
 \end{cases}
\]
On a center-free outgoing edge this gives
\[
 A(T_{\mathrm{next}})\ge
 \begin{cases}
 a,&0\le c\le1/2,\\
 g_c^\vee(a)\ge a,&1/2<c\le1.
 \end{cases}
\]
\item With a center interval $J$ on the handoff edge, every V triangle gives
\[
 A(T_{\mathrm{next}})\ge g_{c,J}^\vee(a).
\]
If $T$ is nonsupercritical, use the sharper branchwise bound
\[
 A(T_{\mathrm{next}})\ge
 \begin{cases}
 \mathcal R_J(1-a),&0\le c\le1/2,\\
 g_{c,J}^\vee(a),&1/2<c\le1.
 \end{cases}
\]
\item If $T$ is strictly supercritical, then necessarily $0\le c<1/2$, and
\[
 B(T)<g_c^{\rm sc}.
\]
Only when its outgoing edge is center-free may this be complemented to
\[
 A(T_{\mathrm{next}})>1-g_c^{\rm sc}.
\]
\end{enumerate}
\end{proposition}

\begin{proof}
Items 1--3 are Lemmas~\ref{lem:reader-raw-capped-transfer} and
\ref{lem:reader-generalized-handoff}.  Item 4 is
\eqref{eq:reader-free-outgoing-envelope}, with the additional center-free
hypothesis used only in \eqref{eq:reader-free-raw-envelope}.
\end{proof}
